% Algebra 2 Lesson 6-3 fast-track packet: Logarithms (enVision A2 2024 TE, Example 1 and Example 5 only)
% Bridge lesson, re-cut 2026-10-31 toward the Topic 6 test. Topic items: 7 = the inverse of y = 4^x
% (Example 1), 9 = solve log_4(5x+9) = 3 (Example 5). Quiz items (Savvas 6-3 Lesson Quiz): 1 = the
% inverse of y = 3^x (Example 1), 3 = log form to exponential form (Example 5 Part B's step, bonus),
% 5 = solve log(2x+4)=2 (Example 5). Cut: Examples 2, 3, 4, 6 (the box below keeps the TE concept
% summary and the common/natural log notation; items 2 and 4 are full Savvas quiz only).
% Shared preamble for the "fast" (hacked-down) Algebra 2 lesson packets, Oct 2026.
% Each packet: % Shared preamble for the "fast" (hacked-down) Algebra 2 lesson packets, Oct 2026.
% Each packet: % Shared preamble for the "fast" (hacked-down) Algebra 2 lesson packets, Oct 2026.
% Each packet: \input{fast-packet-template} then \begin{document} ... \end{document}.
\documentclass[12pt]{article}
\usepackage[margin=0.7in]{geometry}
\usepackage{amsmath,amssymb}
\usepackage{booktabs,tabularx,array}
\usepackage{enumitem}
\usepackage[T1]{fontenc}
\usepackage[utf8]{inputenc}
\usepackage{lmodern}
\usepackage{tikz}
\usepackage{pgfplots}
\pgfplotsset{compat=1.18}
\usepackage{xcolor}
\setlist{nosep,leftmargin=1.4em}
\setlength{\parindent}{0pt}
\setlength{\parskip}{6pt}
\renewcommand{\arraystretch}{1.35}
\newcolumntype{Y}{>{\raggedright\arraybackslash}X}
\pagestyle{empty}

% blanks and workspace
\newcommand{\blank}[1][2.6cm]{\rule[-2pt]{#1}{0.5pt}}
\newcommand{\lines}[1][1]{\par\vspace{2pt}\foreach \i in {1,...,#1}{\rule{0pt}{22pt}\rule{\linewidth}{0.4pt}\par}}
\newcommand{\wspart}[2]{\vspace{4pt}{\large\bfseries #1: #2}\par\vspace{-6pt}\rule{\textwidth}{0.4pt}\par}

% header: \packethead{Lesson 1-3}{Piecewise-Defined Functions}
\newcommand{\packethead}[2]{%
{\LARGE\bfseries #2}\hfill Name \blank[5cm]\\[2pt]
Algebra 2 \quad #1 (fast track) \hfill Date \blank[2.5cm] \quad Period \blank[1.2cm]\par\vspace{2pt}}

% boxed goals: \goals{learning goal}{language goal}
\newcommand{\goals}[2]{%
\fbox{\parbox{\dimexpr\textwidth-2\fboxsep-2\fboxrule}{%
\textbf{Learning goal:} #1\\[2pt]
\textbf{Language goal:} #2}}\par}

% quiz target box: \quiztarget{text}
\newcommand{\quiztarget}[1]{%
\vspace{4pt}\colorbox{gray!12}{\parbox{\dimexpr\textwidth-2\fboxsep}{\textbf{On the quiz:} #1}}\par}

% IXL table row helper: code & skill & which quiz question it feeds
\newcommand{\ixlrow}[3]{\texttt{#1} & #2 & #3 \\}

% one grid style for every graph; override xmin/xmax/... per plot
\pgfplotsset{sheetgrid/.style={
  axis lines=middle, grid=both, grid style={gray!45},
  xmin=-10.5, xmax=10.5, ymin=-10.5, ymax=10.5,
  xtick={-10,-8,...,10}, ytick={-10,-8,...,10},
  x=0.28cm, y=0.28cm,
  ticklabel style={font=\scriptsize, fill=white, inner sep=1pt},
  xlabel={$x$}, ylabel={$y$},
  axis line style={thick}, clip=true}}
% unit grid for small windows
\pgfplotsset{unitgrid/.style={sheetgrid,
  xmin=-5.5, xmax=5.5, ymin=-5.5, ymax=5.5,
  xtick={-5,...,5}, ytick={-5,...,5}, x=0.5cm, y=0.5cm}}
% open and closed endpoints
\tikzset{closed/.style={circle, fill, inner sep=1.6pt},
         open/.style={circle, draw, fill=white, inner sep=1.4pt}}
 then \begin{document} ... \end{document}.
\documentclass[12pt]{article}
\usepackage[margin=0.7in]{geometry}
\usepackage{amsmath,amssymb}
\usepackage{booktabs,tabularx,array}
\usepackage{enumitem}
\usepackage[T1]{fontenc}
\usepackage[utf8]{inputenc}
\usepackage{lmodern}
\usepackage{tikz}
\usepackage{pgfplots}
\pgfplotsset{compat=1.18}
\usepackage{xcolor}
\setlist{nosep,leftmargin=1.4em}
\setlength{\parindent}{0pt}
\setlength{\parskip}{6pt}
\renewcommand{\arraystretch}{1.35}
\newcolumntype{Y}{>{\raggedright\arraybackslash}X}
\pagestyle{empty}

% blanks and workspace
\newcommand{\blank}[1][2.6cm]{\rule[-2pt]{#1}{0.5pt}}
\newcommand{\lines}[1][1]{\par\vspace{2pt}\foreach \i in {1,...,#1}{\rule{0pt}{22pt}\rule{\linewidth}{0.4pt}\par}}
\newcommand{\wspart}[2]{\vspace{4pt}{\large\bfseries #1: #2}\par\vspace{-6pt}\rule{\textwidth}{0.4pt}\par}

% header: \packethead{Lesson 1-3}{Piecewise-Defined Functions}
\newcommand{\packethead}[2]{%
{\LARGE\bfseries #2}\hfill Name \blank[5cm]\\[2pt]
Algebra 2 \quad #1 (fast track) \hfill Date \blank[2.5cm] \quad Period \blank[1.2cm]\par\vspace{2pt}}

% boxed goals: \goals{learning goal}{language goal}
\newcommand{\goals}[2]{%
\fbox{\parbox{\dimexpr\textwidth-2\fboxsep-2\fboxrule}{%
\textbf{Learning goal:} #1\\[2pt]
\textbf{Language goal:} #2}}\par}

% quiz target box: \quiztarget{text}
\newcommand{\quiztarget}[1]{%
\vspace{4pt}\colorbox{gray!12}{\parbox{\dimexpr\textwidth-2\fboxsep}{\textbf{On the quiz:} #1}}\par}

% IXL table row helper: code & skill & which quiz question it feeds
\newcommand{\ixlrow}[3]{\texttt{#1} & #2 & #3 \\}

% one grid style for every graph; override xmin/xmax/... per plot
\pgfplotsset{sheetgrid/.style={
  axis lines=middle, grid=both, grid style={gray!45},
  xmin=-10.5, xmax=10.5, ymin=-10.5, ymax=10.5,
  xtick={-10,-8,...,10}, ytick={-10,-8,...,10},
  x=0.28cm, y=0.28cm,
  ticklabel style={font=\scriptsize, fill=white, inner sep=1pt},
  xlabel={$x$}, ylabel={$y$},
  axis line style={thick}, clip=true}}
% unit grid for small windows
\pgfplotsset{unitgrid/.style={sheetgrid,
  xmin=-5.5, xmax=5.5, ymin=-5.5, ymax=5.5,
  xtick={-5,...,5}, ytick={-5,...,5}, x=0.5cm, y=0.5cm}}
% open and closed endpoints
\tikzset{closed/.style={circle, fill, inner sep=1.6pt},
         open/.style={circle, draw, fill=white, inner sep=1.4pt}}
 then \begin{document} ... \end{document}.
\documentclass[12pt]{article}
\usepackage[margin=0.7in]{geometry}
\usepackage{amsmath,amssymb}
\usepackage{booktabs,tabularx,array}
\usepackage{enumitem}
\usepackage[T1]{fontenc}
\usepackage[utf8]{inputenc}
\usepackage{lmodern}
\usepackage{tikz}
\usepackage{pgfplots}
\pgfplotsset{compat=1.18}
\usepackage{xcolor}
\setlist{nosep,leftmargin=1.4em}
\setlength{\parindent}{0pt}
\setlength{\parskip}{6pt}
\renewcommand{\arraystretch}{1.35}
\newcolumntype{Y}{>{\raggedright\arraybackslash}X}
\pagestyle{empty}

% blanks and workspace
\newcommand{\blank}[1][2.6cm]{\rule[-2pt]{#1}{0.5pt}}
\newcommand{\lines}[1][1]{\par\vspace{2pt}\foreach \i in {1,...,#1}{\rule{0pt}{22pt}\rule{\linewidth}{0.4pt}\par}}
\newcommand{\wspart}[2]{\vspace{4pt}{\large\bfseries #1: #2}\par\vspace{-6pt}\rule{\textwidth}{0.4pt}\par}

% header: \packethead{Lesson 1-3}{Piecewise-Defined Functions}
\newcommand{\packethead}[2]{%
{\LARGE\bfseries #2}\hfill Name \blank[5cm]\\[2pt]
Algebra 2 \quad #1 (fast track) \hfill Date \blank[2.5cm] \quad Period \blank[1.2cm]\par\vspace{2pt}}

% boxed goals: \goals{learning goal}{language goal}
\newcommand{\goals}[2]{%
\fbox{\parbox{\dimexpr\textwidth-2\fboxsep-2\fboxrule}{%
\textbf{Learning goal:} #1\\[2pt]
\textbf{Language goal:} #2}}\par}

% quiz target box: \quiztarget{text}
\newcommand{\quiztarget}[1]{%
\vspace{4pt}\colorbox{gray!12}{\parbox{\dimexpr\textwidth-2\fboxsep}{\textbf{On the quiz:} #1}}\par}

% IXL table row helper: code & skill & which quiz question it feeds
\newcommand{\ixlrow}[3]{\texttt{#1} & #2 & #3 \\}

% one grid style for every graph; override xmin/xmax/... per plot
\pgfplotsset{sheetgrid/.style={
  axis lines=middle, grid=both, grid style={gray!45},
  xmin=-10.5, xmax=10.5, ymin=-10.5, ymax=10.5,
  xtick={-10,-8,...,10}, ytick={-10,-8,...,10},
  x=0.28cm, y=0.28cm,
  ticklabel style={font=\scriptsize, fill=white, inner sep=1pt},
  xlabel={$x$}, ylabel={$y$},
  axis line style={thick}, clip=true}}
% unit grid for small windows
\pgfplotsset{unitgrid/.style={sheetgrid,
  xmin=-5.5, xmax=5.5, ymin=-5.5, ymax=5.5,
  xtick={-5,...,5}, ytick={-5,...,5}, x=0.5cm, y=0.5cm}}
% open and closed endpoints
\tikzset{closed/.style={circle, fill, inner sep=1.6pt},
         open/.style={circle, draw, fill=white, inner sep=1.4pt}}

\begin{document}

\packethead{Lesson 6-3}{Logarithms}
\goals{I can write the inverse of an exponential function as a logarithm, and solve an equation by switching between logarithmic form and exponential form.}{I can explain how I found an inverse or solved an equation using the words base, exponent, argument, inverse, and logarithmic form.}

\quiztarget{(1) Name the inverse of an exponential function: Example 1, Practice. (2) Estimate $\ln$ of a number (full Savvas quiz only). (3) Rewrite a logarithmic equation in exponential form: Example 5 Part B's first step (bonus). (4) Select the logarithms evaluated correctly (full Savvas quiz only). (5) Solve an equation with a common logarithm in it: Example 5.}
\topictarget{6}{Question 7 asks which function is the inverse of an exponential function (Example 1). Question 9 asks you to solve an equation with a base-4 logarithm in it (Example 5).}

\vspace{4pt}
\fbox{\parbox{\dimexpr\textwidth-2\fboxsep-2\fboxrule}{%
\textbf{Exponential and logarithmic forms.} A logarithm \emph{is} an exponent.\par
\renewcommand{\arraystretch}{1.15}%
\begin{tabularx}{\linewidth}{@{}l Y Y@{}}
 & \textbf{Exponential form} & \textbf{Logarithmic form} \\
Algebra & $b^x=y$ & $\log_b y=x$ \\
Words & the \textbf{base} raised to the \textbf{exponent} is the result & the log, base $b$, of the result (the \textbf{argument}) is the exponent \\
Numbers & $3^4=81$ & $\log_3 81=4$ \\
\end{tabularx}\renewcommand{\arraystretch}{1.35}\\[2pt]
$\log x$ means base 10 (\textbf{common logarithm}); $\ln x$ means base $e$ (\textbf{natural logarithm}). Spanish: logaritmo.}}

\wspart{Example 1}{Understand Logarithms}
Solve $2^x=8$. The operation in $2^x=8$ is exponentiation. To solve this equation, you need an inverse operation that answers the question, ``To what exponent would you raise the base 2 to get 8?''

\begin{minipage}[t]{0.48\textwidth}
\textbf{Compare:} $2x=8$. Division is the inverse of multiplication, so you can divide both sides by 2:
$\frac{2x}{2}=\frac82$, \ $x=4$.
\end{minipage}\hfill
\begin{minipage}[t]{0.48\textwidth}
\textbf{Now:} $2^x=8$. A logarithm is the inverse of exponentiation, so you can write
$\log_2 8=x$ (read ``log base 2 of 8''). \\
$\log_2 8=3$ \ because \ $2^3=8$.
\end{minipage}

\vspace{4pt}
\textbf{Definition.} $\log_b x=y$ \ if and only if \ $b^y=x$, \ for $b>0$, $b\neq1$, and $x>0$.\\
The logarithmic function $y=\log_b x$ is the \textbf{inverse} of the exponential function $y=b^x$.

\textbf{Say it:} To solve $2^x=8$, I ask what exponent on the base \blank[1.2cm] gives the argument \blank[1.2cm]. In logarithmic form, $\log_2 8=$ \blank[1.2cm] because \blank[2cm]. The inverse of $y=b^x$ is $y=$ \blank[2.4cm].

\wspart{Try It 1}{Understand Logarithms}
Write the logarithmic form of $y=8^x$.

Base \blank[1cm] \ exponent \blank[1cm] \ result \blank[1cm] \qquad Logarithmic form: \blank[4cm]
\lines[1]

\newpage

% R5: structural parallel of Savvas quiz item 1 ("Which function is the inverse of the exponential
% function y = 3^x?"; also Topic 6 item 7). Try It 1 writes the log form x = log_8 y but never names
% the inverse function. Different bases below.
\wspart{Practice}{Name the inverse}
Write the inverse of each exponential function as a logarithmic function $y=\log_b x$.

\begin{minipage}[t]{0.31\textwidth}
\textbf{a.} $y=5^x$
\lines[2]
Inverse: $y=$ \blank[2cm]
\end{minipage}\hfill
\begin{minipage}[t]{0.31\textwidth}
\textbf{b.} $y=10^x$
\lines[2]
Inverse: $y=$ \blank[2cm]
\end{minipage}\hfill
\begin{minipage}[t]{0.31\textwidth}
\textbf{c.} $y=\left(\tfrac12\right)^x$
\lines[2]
Inverse: $y=$ \blank[2cm]
\end{minipage}

\vspace{4pt}
In part b, how do you write the inverse without writing the base? (Look at the box on page 1.)\lines[1]

\wspart{Example 5}{Solve Equations With Logarithms}
What is the solution to each equation? Round to the nearest thousandth.

\begin{minipage}[t]{0.48\textwidth}
\textbf{A.} $25=10^{x-1}$

The variable is in the \textbf{exponent}. Convert to logarithmic form (base 10, exponent \mbox{$x-1$}, result 25):
\begin{enumerate}[label=\arabic*.]
\item $25=10^{x-1}$
\item $\log25=x-1$
\item $1+\log25=x$
\item $2.398\approx x$
\end{enumerate}
\textbf{Solution:} $x\approx2.398$.

\textbf{Common error:} 10 is not a coefficient, it is a base. You cannot divide both sides by 10 and then add 1.
\end{minipage}\hfill
\begin{minipage}[t]{0.48\textwidth}
\textbf{B.} $\ln(2x+3)=4$

The variable is in the \textbf{argument}. Convert to exponential form (base $e$, exponent 4, result \mbox{$2x+3$}):
\begin{enumerate}[label=\arabic*.]
\item $\ln(2x+3)=4$
\item $2x+3=e^4$
\item $2x+3\approx54.598$
\item $2x\approx51.598$
\item $x\approx25.799$
\end{enumerate}
\textbf{Solution:} $x\approx25.799$.

\textbf{Check:} $\ln\bigl(2(25.799)+3\bigr)=\ln 54.598\approx4$. \checkmark
\end{minipage}

\vspace{4pt}
\textbf{Say it:} In Part B, I rewrote the logarithmic form $\ln(2x+3)=4$ in exponential form. The base is \blank[1.2cm], the exponent is \blank[1.2cm], and the argument $2x+3$ equals \blank[1.5cm]. Then I \dots \lines[2]

\newpage

\wspart{Try It 5}{Solve Equations With Logarithms}
Solve each equation. Round to the nearest thousandth.

\begin{minipage}[t]{0.48\textwidth}
\textbf{a.} $\log(3x-2)=2$

Base \blank[1cm] \ exponent \blank[1cm] \ argument \blank[1.6cm]

Exponential form: \blank[3.4cm]
\lines[4]
$x=$ \blank[2.4cm]
\end{minipage}\hfill
\begin{minipage}[t]{0.48\textwidth}
\textbf{b.} $e^{x+2}=8$

Base \blank[1cm] \ exponent \blank[1cm] \ result \blank[1.6cm]

Logarithmic form: \blank[3.4cm]
\lines[4]
$x\approx$ \blank[2.4cm]
\end{minipage}

\vspace{4pt}
Why did you switch to exponential form in part a but to logarithmic form in part b? Use the words \emph{exponent} and \emph{argument}.\lines[2]

\wspart{IXL}{Practice for the quiz}
Go to IXL, Algebra 2. Type the code in the search box. Do the starred skills first.

\vspace{2pt}
\begin{tabularx}{\textwidth}{l Y Y}
\toprule
\textbf{Code} & \textbf{Skill} & \textbf{Feeds} \\
\midrule
\ixlrow{8RK $\bigstar$}{Convert between exponential and logarithmic form: all bases}{Quiz 1 and 3; the first step of Quiz 5 (Examples 1 and 5, Practice)}
\ixlrow{BXU $\bigstar$}{Solve logarithmic equations I}{Quiz 5 (Example 5, Try It 5a)}
\ixlrow{GBR}{Evaluate logarithms (optional)}{Quiz 4 (full Savvas quiz only)}
\ixlrow{TDF}{Evaluate logarithms using a calculator (optional)}{Quiz 2 (full Savvas quiz only); the calculator step in Example 5}
\bottomrule
\end{tabularx}

\vspace{4pt}
\textbf{SmartScore 70 or above on each starred ($\bigstar$) skill replaces your quiz score.}

\newpage

% DOK 3: TE Practice & Problem Solving 17 (Understand, Use Structure). The item analysis files every
% DOK 3 item under a cut Example (17 under Example 3; 13, 15, 18 under Example 4; 57 under Example 6),
% so R3 names 17 as the self-paced extension. The box and Example 1's definition are enough to do it.
\wspart{Extension (self-paced)}{Use Structure}
\fbox{\parbox{\dimexpr\textwidth-2\fboxsep-2\fboxrule}{%
A student says that $\log_3\!\left(\frac1{27}\right)$ simplifies to $-3$. Is the student correct? Explain.

\vspace{6pt}
The student is \blank[2.4cm] because the base \blank[1.2cm] raised to the exponent \blank[1.2cm] gives the argument \blank[1.5cm].
\lines[3]\vspace{6pt}}}

\vspace{8pt}
\textbf{Before the quiz, finish these sentences.}

A logarithm is an \blank[3cm]: $\log_b x$ is the exponent you put on the base \blank[1.2cm] to get the argument \blank[1.2cm].

The \blank[2.2cm] of the exponential function $y=b^x$ is the logarithmic function $y=$ \blank[2.4cm].

When the variable is in the exponent, rewrite the equation in \blank[3.4cm] form.

To solve $\log(\text{expression})=c$, rewrite it in \blank[4cm] form: $10^{c}=$ \blank[3cm].

\end{document}
